%===============================================================================
% APPENDIX G: GLOSSARY
%===============================================================================
\chapter{Glossary}
\label{app:glossary}

\begin{description}
\item[aBFT (Asynchronous Byzantine Fault Tolerance)] A consensus property guaranteeing safety and liveness in asynchronous networks even when up to 1/3 of nodes are malicious.

\item[Blockchain] A distributed ledger technology where records (blocks) are cryptographically linked and immutable.

\item[Client-Side Trust] Security assumption that the user's browser/device executes code as intended, without malicious modification.

\item[Consensus Timestamp] The agreed-upon time assigned to a transaction by the network, ensuring fair ordering.

\item[CORS (Cross-Origin Resource Sharing)] Browser security mechanism controlling which origins can access resources.

\item[DApp (Decentralized Application)] An application that runs on a decentralized network rather than a single server.

\item[DLT (Distributed Ledger Technology)] The technological infrastructure underlying blockchain systems.

\item[Gas] Computational fee unit on Ethereum and similar platforms; not used in Hedera HCS.

\item[HCS (Hedera Consensus Service)] Hedera's topic-based messaging service providing fair ordering and consensus timestamps.

\item[HBAR] Native cryptocurrency of the Hedera network, used for transaction fees.

\item[Hashgraph] Hedera's consensus algorithm using virtual voting and gossip protocol for aBFT.

\item[HCS Topic] A named channel on Hedera for submitting ordered, timestamped messages.

\item[JWT (JSON Web Token)] Compact, URL-safe token format for transmitting claims between parties.

\item[KYC (Know Your Customer)] Identity verification process required by financial regulations.

\item[Mainnet] The production, live version of a blockchain network with real economic value.

\item[Mirror Node] Read-only Hedera node providing historical query access to consensus data.

\item[Nonce] A number used once, typically for preventing replay attacks.

\item[Optimistic Locking] Concurrency control strategy assuming conflicts are rare, checking at commit time.

\item[Permissioned Ledger] Blockchain where validator identities are known and approved (e.g., Hyperledger Fabric).

\item[Proof of Stake (PoS)] Consensus mechanism where validators stake tokens as collateral.

\item[Proof of Work (PoW)] Consensus mechanism requiring computational puzzle solving.

\item[Receipt-freeness] Voting property preventing voters from proving how they voted to a coercer.

\item[Rollback] Reverting a database change when a subsequent operation fails.

\item[Smart Contract] Self-executing code on a blockchain that enforces predefined rules.

\item[Testnet] A test version of a blockchain network with no real economic value.

\item[TPS (Transactions Per Second)] Measure of blockchain throughput capacity.

\item[Virtual Voting] Hashgraph's method of calculating consensus without actual voting messages.

\item[Web2] Traditional centralized web architecture (servers, databases, client-server).

\item[Web3] Decentralized web architecture using blockchain and peer-to-peer protocols.

\item[Web2.5] Hybrid architecture combining Web2 identity/UX with Web3 consensus/verifiability.

\item[ZK-Proof (Zero-Knowledge Proof)] Cryptographic proof that a statement is true without revealing secret inputs.
\end{description}